\section{Proof details}\label{app:proofs}

\subsection{Encoding lengths for Theorem~\ref{thm:polytope}}\label{app:polytope}

Multiply each row of $Ax\le b$ by the product of the denominators of its
entries, and multiply $q$ by the product $\sigma$ of the denominators of its
coefficients. Positive row scalings and $\sigma>0$ change neither $P$, nor
the nonsingularity of $K_S$, nor the $x$-part of its solution. Afterwards
all row entries are integers of absolute value at most $2^{\alpha}$ and all
entries of $\sigma H$ and $\sigma c$ are integers of absolute value at most
$2^{\eta}$, where $\alpha$ is at most the largest row encoding length and
$\eta$ at most the encoding length of $q$ plus one.

Let $\abs S=k\le d$ and $n=d+k$. Scale the last $k$ rows and the last $k$
columns of the integer matrix $[K_S\mid(-\sigma c,b_S)]$ by $2^{-\alpha}$. In
the scaled matrix, the first $d$ rows have entries of absolute value at
most $2^\eta$ and the others at most $1$. Hadamard's inequality applied to
any square submatrix of order $s\le n$ and the undoing of the scaling give
\[
\abs{\det M}\le n^{n/2}2^{d\eta+2k\alpha}\le B:=(2d)^d2^{d(\eta+2\alpha)}
\]
for every square submatrix $M$ of $[K_S\mid(-\sigma c,b_S)]$. By Cramer's
rule a candidate has coordinates $x_j=N_j/\Delta$ with
$\abs{N_j},\abs\Delta\le B$. Writing $\sigma q=\tilde c_0+\sum_j\tilde c_jx_j+\sum_{i\le j}\tilde Q_{ij}x_ix_j$
with integer coefficients bounded by $2^\eta$,
\[
\sigma q(x)=\frac{\Delta^2\tilde c_0+\Delta\sum_j\tilde c_jN_j+\sum_{i\le j}\tilde Q_{ij}N_iN_j}{\Delta^2},
\]
a numerator with at most $(d+1)^2$ terms, each bounded by $2^\eta B^2$. So
$\min_Pq$ is a fraction whose numerator and denominator are bounded by
$(d+1)^22^{\eta}B^2$ and $2^{\eta}B^2$. Hence every candidate and
the minimum have bit length $O(d(\eta+\alpha+\log d))$, for every $d$.

Bareiss's fraction-free elimination keeps every intermediate number equal
to a minor of the matrix above, and the products formed before each exact
division are bounded by $2B^2$. Feasibility tests and value comparisons
are integer comparisons of numbers bounded by $(d+1)2^\alpha B$ and
$(d+1)^22^\eta B^4$. The total work is $O(N(m,d)(d^3+md))$ arithmetic
operations on integers of bit length $O(d(\eta+\alpha+\log d))$. For fixed
$d$, $N(m,d)\le(m+1)^d$, the number of arithmetic operations does not
depend on the sizes of the numbers, and the algorithm is strongly
polynomial.

\subsection{Lower bounds for Theorem~\ref{thm:interleave}}\label{app:interleave}

Let $\nu_A$ and $\nu_C$ be probability measures on $[\ell,u]$ with equal
moments of orders one and two. For any polynomial $p$ of degree at most
two, $\int\dist(y,A)^2d\nu_A+\int\dist(y,C)^2d\nu_C$ equals
$\int(\dist(y,A)^2-p)\,d\nu_A+\int(\dist(y,C)^2+p)\,d\nu_C$, so it suffices to
find $p$ with $\dist(\cdot,A)^2-p\ge\kappa_A$ and $\dist(\cdot,C)^2+p\ge\kappa_C$
on $\R$ and $\kappa_A+\kappa_C=\delta^2/2$. Because $\dist(y,S)^2=\min_{s\in S}(y-s)^2$,
each inequality can be checked one point $s$ at a time. We use
$\inf_y[\alpha(y-s)^2+\beta(y-m)^2]=\alpha\beta(s-m)^2/(\alpha+\beta)$ for
$\alpha+\beta>0$.

\emph{Separated sets.} By symmetry let $\max A<\min C$, put $\alpha_0=\max A$,
$\gamma_0=\min C$, $\delta=\gamma_0-\alpha_0$, $m_0=(\alpha_0+\gamma_0)/2$ and
$p(y)=\delta(y-m_0)$. For $s\in A$,
$\inf_y[(y-s)^2-\delta(y-m_0)]=-\delta^2/4-\delta(s-m_0)\ge\delta^2/4$, and for
$t\in C$, $\inf_y[(y-t)^2+\delta(y-m_0)]=-\delta^2/4+\delta(t-m_0)\ge\delta^2/4$.

\emph{Nested sets.} By symmetry let $C$ lie strictly between the two
points $a_-<a_+$ of $A$. Put $m=(\min C+\max C)/2$, $R_0=(\max C-\min C)/2$
and $r=\delta+R_0$; then $\abs{s-m}\ge r$ for $s\in A$ and $\abs{t-m}\le R_0$ for
$t\in C$. If $R_0>0$, let $\gamma=(r-R_0)/(r+R_0)\in(0,1)$ and
$p(y)=-\gamma(y-m)^2$. Then
$\inf_y[(y-s)^2+\gamma(y-m)^2]=\gamma(s-m)^2/(1+\gamma)\ge\gamma r^2/(1+\gamma)$ and
$\inf_y[(y-t)^2-\gamma(y-m)^2]=-\gamma(t-m)^2/(1-\gamma)\ge-\gamma R_0^2/(1-\gamma)$.
With $1+\gamma=2r/(r+R_0)$ and $1-\gamma=2R_0/(r+R_0)$ the two bounds add up to
$(r-R_0)^2/2=\delta^2/2$. If $R_0=0$, take $p(y)=-(y-m)^2$; then
$\dist(\cdot,C)^2+p\equiv0$ and $\dist(\cdot,A)^2-p\ge\delta^2/2$.

\subsection{A separating polynomial for Theorem~\ref{thm:alternation}}\label{app:alternation}

Let $A$ and $C$ be disjoint, nonempty and compact, let $T=A\cup C$, and
label $t\in T$ by $+1$ if $t\in A$ and $-1$ if $t\in C$. Let
$\varepsilon=\dist(A,C)>0$. Starting from $s_0=\min T$, let $s_{i+1}$ be the
smallest point of $T$ larger than $s_i$ with the opposite label, as long
as one exists. Such a point exceeds $s_i$ by at least $\varepsilon$, so
the process stops at some $s_r$, and every point of $T$ in $[s_i,s_{i+1})$
has the label of $s_i$. Every alternating chain visits these runs in
increasing order, so $\mathrm{alt}(A,C)=r$. Let $e_i$ be the largest point
of $T$ in $[s_{i-1},s_i)$ and choose $\rho_i\in(e_i,s_i)$. Then
$p(t)=\lambda(s_0)\prod_{i=1}^r(\rho_i-t)$ has degree $r$, and a point of the
$i$-th run lies between $\rho_i$ and $\rho_{i+1}$, so
$\operatorname{sign}p(t)$ equals the label of $t$: $p>0$ on $A$ and $p<0$ on
$C$. If $r\le k$, then $\int p\,d\alpha=\int p\,d\gamma$ for any probability
measures $\alpha$ on $A$ and $\gamma$ on $C$ with equal moments up to order
$k$, which is impossible.

\subsection{Dense moment matrices for the path family}\label{app:sdp}

We use the notation of Theorem~\ref{thm:interleave} and assume that $A$
and $C$ are disjoint. Replacing $x$ by $1-x$ exchanges $a_1$ and $a_2$ and
maps the four McCormick inequalities of the product $xz$ to each other,
and similarly for $z$. We may therefore assume $a_1<a_2$ and $c_1<c_2$.

\begin{proposition}\label{prop:sdp}
Let $v\in R$ be a point at which the linearization of $D$ vanishes.
There is no value of $p_{xz}$ for which the moment matrix
\[
M=\begin{pmatrix}1&m_x&m_y&m_z\\ m_x&s_x&p_{xy}&p_{xz}\\ m_y&p_{xy}&s_y&p_{yz}\\ m_z&p_{xz}&p_{yz}&s_z\end{pmatrix}
\]
is positive semidefinite and both $p_{xz}\le m_x$ and $p_{xz}\le m_z$ hold.
\end{proposition}

\begin{proof}
By the proof of Theorem~\ref{thm:interleave}, the linearization of $D_A$
vanishes at the $(x,y)$ part of $v$, so that part is the moment vector of
a measure supported on the zero set of $D_A$. This zero set is
$\{(0,a_1),(1,a_2)\}$: at a point with $0<x<1$, concavity or affinity in
$x$ gives $D_A(x,y)\ge(1-x)(y-a_1)^2+x(y-a_2)^2>0$. Hence $x=\alpha+\beta y$
almost surely on this pair, with $\beta=1/(a_2-a_1)$ and
$\alpha=-a_1\beta$, and the vector $w=(\alpha,-1,\beta,0)$ satisfies
$w\T Mw=0$. If $M\succeq0$, then $Mw=0$; the last row gives
$p_{xz}=\alpha m_z+\beta p_{yz}$. In the same way $z=(y-c_1)/(c_2-c_1)$
almost surely on the $(y,z)$ pair, which gives
$m_z=(m_y-c_1)/(c_2-c_1)$ and $p_{yz}=(s_y-c_1m_y)/(c_2-c_1)$. Combining,
\[
p_{xz}=\frac{s_y-(a_1+c_1)m_y+a_1c_1}{(a_2-a_1)(c_2-c_1)}.
\]
The numerator is the expectation of $(y-a_1)(y-c_1)$ under any measure
with moments $(m_y,s_y)$. Evaluating it with the $(x,y)$ measure, which
puts mass $m_x$ on $y=a_2$, gives $p_{xz}=m_x(a_2-c_1)/(c_2-c_1)$;
evaluating it with the $(y,z)$ measure gives
$p_{xz}=m_z(c_2-a_1)/(a_2-a_1)$.

By Theorem~\ref{thm:interleave}, $A$ and $C$ interleave, and since the
chords meet at a point interior to both, $0<m_x<1$ and $0<m_z<1$. If
$a_1<c_1<a_2<c_2$, then $(c_2-a_1)/(a_2-a_1)>1$ and $p_{xz}>m_z$. If
$c_1<a_1<c_2<a_2$, then $(a_2-c_1)/(c_2-c_1)>1$ and $p_{xz}>m_x$.
\end{proof}

The positive semidefinite completion in this proof is unique, so the
dense moment matrix alone does not exclude the zero-value point, and the
McCormick inequalities alone admit every $p_{xz}$ between
$\max\{0,m_x+m_z-1\}$ and $\min\{m_x,m_z\}$. Together they do, and in fact
they prove $D\ge\delta^2/2$: in~\eqref{eq:sdp-identity}, write
$u^2+v^2=\tfrac12(u+v)^2+\tfrac12(u-v)^2$; the function
$g=\tfrac12(u-v)^2+\tfrac12(a_2-a_1)^2x(1-x)+\tfrac12(c_2-c_1)^2z(1-z)$ has no
$x^2$ or $z^2$ terms, so it is bilinear in $(x,z)$ and equals its
interpolation $\sum_{ij}g(i,j)\ell_{ij}$ from the four vertices, where
$g(i,j)=\tfrac12(c_{j+1}-a_{i+1})^2$; and $\sum_{ij}\ell_{ij}=1$.

\subsection{Sizes, a lower bound and a depth-two example for Theorem~\ref{thm:star}}\label{app:star}

For a rational $r=p/q$ in lowest terms let $\eta(r)=\log_2\max\{\abs p,q\}$.
Then $\eta(rs),\eta(r/s)\le\eta(r)+\eta(s)$ and
$\eta(r_1+\dots+r_N)\le\log_2N+\sum_j\eta(r_j)$. Let $\sigma_i$ bound $\eta$
over the data of leaf $i$ (its coefficients, bounds and rows), $\sigma_0$
over the center data, and $\sigma=\max_i\sigma_i$.

\emph{Sizes.} The coefficients of the affine functions defining $L_i$ and
$U_i$ have $\eta\le2\sigma_i$, and $v_i$ has $\eta\le2\sigma_i+1$. A breakpoint is
the intersection of two such functions or the zero of the affine factor
in~\eqref{eq:endpoint-factor}; in both cases $\eta\le12\sigma+4$, independently
of $k$. On a piece, the contribution of leaf $i$ to $V$ is a quadratic in
$y$ whose coefficients have $\eta\le11\sigma_i+7$, so the coefficients of $V$,
which are sums of $k+1$ such terms, have $\eta\le C:=\log_2(k+1)+\sigma_0+\sum_i(11\sigma_i+7)$,
which is linear in the input length. Evaluating at a breakpoint or at a
stationary point $-c_1/(2c_2)$ gives a minimum with $\eta\le4C+36\sigma+14$.
Because the sweep updates the coefficients exactly, every running sum
equals the coefficient of the current piece and obeys the same bound. The
algorithm is therefore strongly polynomial.

\emph{Lower bound.} Ben-Or's theorem states that an algebraic computation
tree deciding membership in a set $W\subseteq\R^n$ with $N$ connected
components has depth $\Omega(\log N-n)$ \citep{BenOr1983}. For $r\ge1$ let
$W=\{z\in(0,2r)^r:\abs{z_i-z_j}\ge1\ \text{for}\ i\ne j\}$, which has one
component for each of the $r!$ orderings. For $z\in(0,2r)^r$ build a star
with center $y\in[0,2r]$ and, for each $j$, three leaves: $x\in[0,1]$ with
objective $(y-z_j+1)x$; $x\in[0,1]$ with objective $(z_j+1-y)x$; and
$x\in[0,2r+1]$ with the row $x\ge y-z_j$ and objective $2x$; add the center
terms $-ry+\sum_jz_j-r$. With $t=y-z_j$, the identity
$\min\{0,t+1\}+\min\{0,1-t\}+2\max\{0,t\}-t-1=-\max\{0,1-\abs t\}$ shows that
$V(y)=-\sum_j\max\{0,1-\abs{y-z_j}\}$, a sum of negative tents of width two.
If $z\in W$, at most two tents are positive at any $y$ and their sum is at
most one, so $\min V\ge-1$; if $\abs{z_i-z_j}<1$, then $V<-1$ at the midpoint of
$z_i$ and $z_j$. The instance has $3r$ leaves and $r$ leaf rows and is
computed from $z$ with $O(r)$ operations, so deciding $\min V\ge-1$ needs
depth $\Omega(r\log r)$.

\emph{Depth two.} For $f=(x_2-\tfrac18)x_1+x_2^2-x_2+x_2x_3+x_3^2-x_3$ on
$[0,1]^3$, minimizing over $x_1$ first gives
$\min\{0,x_2-\tfrac18\}+x_2^2+(t-1)x_2$ with $t=x_3$. On $[0,\tfrac18]$ its
minimum is $-\tfrac18$ at $x_2=0$; on $[\tfrac18,1]$ it is $-\tfrac14(1-t)^2$ at
$x_2=(1-t)/2$ when $t\le\tfrac34$, and larger than $-\tfrac18$ otherwise. Hence the
conditional value is $\min\{-\tfrac18,-\tfrac14(1-t)^2\}$, whose two pieces meet
at $t=1-\sqrt2/2$, a root of $2t^2-4t+1$. Adding $t^2-t$ and minimizing over
$t\in[0,1]$ gives $\min f=-\tfrac38$ at $(1,0,\tfrac12)$. Rooting the path at
$x_2$ instead makes it a star, for which the partition is rational.


\subsection{An oracle-polynomial variant of Theorem~\ref{thm:weak-separation}}\label{app:ellipsoid}

Assume $R>0$, put $\eta=(\epsilon-\delta)/2$, $\rho=\min\{\eta/R,1\}$ and
$\nu=\rho^k/2$, and run the central-cut ellipsoid method of
\citet[Theorem~3.2.1]{GrotschelLovaszSchrijver1988} in $\R^k$ with outer
radius $k$ and volume bound $\nu$, answering a query $y$ as follows. If
$y\notin C_J$, return the coordinate vector of a violated bound. Otherwise
call the oracle at $y$, record $w_y=F(p_y)-\bar z$, keep the largest
violation $\mathrm{LB}=\max\{0,\max_y(L_y-y\T\bar z)\}$ found so far together
with its cut, stop if $w_y=0$, and return $-w_y/\norm{w_y}_\infty$. After
$O(k^2\log(kR/(\epsilon-\delta)))$ calls, solve the master LP~\eqref{eq:master}
over all recorded points. We claim that its value satisfies
$\mathrm{UB}\le\mathrm{LB}+\epsilon$, so that either a cut is available
($\mathrm{LB}>0$) or the master LP gives a within-tolerance certificate.

Let $\mathrm{LB}^\ast$ be the final value of LB and let
$K^\ast=\{c\in C_J:\ w\T c\ge\mathrm{LB}^\ast+\delta+\eta\ \text{for all recorded }w\}$.
Every answer given during the run is a valid separation answer for $K^\ast$:
coordinate answers separate $y$ from $C_J\supseteq K^\ast$, and for a
recorded $w_y$ and $c\in K^\ast$ we have
$w_y\T c\ge\mathrm{LB}^\ast+\delta+\eta>(L_y-y\T\bar z)+\delta\ge w_y\T y$.
Since the method never receives a membership answer, it ends with an
ellipsoid of volume at most $\nu$ that contains $K^\ast$. If
$\mathrm{UB}>\mathrm{LB}^\ast+\epsilon$ at an optimal $\bar c$, then every
$c\in C_J$ with $\norm{c-\bar c}_\infty\le\eta/R$ satisfies
$w\T c\ge\mathrm{UB}-\eta>\mathrm{LB}^\ast+\delta+\eta$ for all recorded $w$, so
$K^\ast$ contains a box of side $\rho$ and has volume at least $\rho^k>\nu$,
a contradiction. The queries have polynomial encoding length, and the
final certificate is checked independently of the ellipsoid computation.

